% Core reference vocabulary; curated clinical classification of genetic and chromosomal diseases.
\small
\noindent\textit{Scope}: This reference-oriented list organizes human genetic diseases by mode of inheritance, cytogenetic classification, and molecular pathophysiological mechanism. Genetic diseases result from alterations in DNA sequences, ranging from single-nucleotide variants to gross chromosomal rearrangements and non-Mendelian mechanisms, providing a focused reference complementing Common Diseases and Rare Diseases.

\medskip
\begin{description}
\item[Autosomal Dominant Diseases] Huntington Disease; Marfan Syndrome; Neurofibromatosis Type 1 (von Recklinghausen Disease); Neurofibromatosis Type 2; Familial Hypercholesterolemia; Polycystic Kidney Disease (Autosomal Dominant, ADPKD); Achondroplasia; Tuberous Sclerosis Complex; Hereditary Spherocytosis; Von Hippel–Lindau Syndrome; Li–Fraumeni Syndrome; Multiple Endocrine Neoplasia (MEN 1 and 2); Myotonic Dystrophy (Type 1 and 2); Hereditary Non-Polyposis Colorectal Cancer (Lynch Syndrome); Familial Adenomatous Polyposis (FAP); Osteogenesis Imperfecta (Type I); Ehlers–Danlos Syndrome (Hypermobile and Classic Types).

\item[Autosomal Recessive Diseases] Cystic Fibrosis; Sickle-Cell Disease; Beta-Thalassemia; Alpha-Thalassemia; Phenylketonuria (PKU); Gaucher Disease; Tay–Sachs Disease; Niemann–Pick Disease; Wilson Disease; Hereditary Hemochromatosis; Alpha-1 Antitrypsin Deficiency; Spinal Muscular Atrophy (SMA); Friedreich Ataxia; Alkaptonuria; Homocystinuria; Maple Syrup Urine Disease; Congenital Adrenal Hyperplasia; Galactosemia; Glycogen Storage Diseases (von Gierke, Pompe, Cori); Ataxia–Telangiectasia; Fanconi Anemia; Xeroderma Pigmentosum; Bloom Syndrome; Ellis–van Creveld Syndrome; Kartagener Syndrome.

\item[X-Linked Recessive Diseases] Duchenne Muscular Dystrophy; Becker Muscular Dystrophy; Hemophilia A (Factor VIII Deficiency); Hemophilia B (Factor IX Deficiency); Glucose-6-Phosphate Dehydrogenase (G6PD) Deficiency; Red-Green Color Blindness; Fragile X Syndrome; Lesch–Nyhan Syndrome; Wiskott–Aldrich Syndrome; X-Linked Agammaglobulinemia (Bruton Tyrosine Kinase Deficiency); X-Linked Adrenoleukodystrophy; Hunter Syndrome (MPS II); Fabry Disease; Menkes Disease; Alport Syndrome (X-Linked Type); Chronic Granulomatous Disease (X-Linked Form).

\item[X-Linked Dominant Diseases] Rett Syndrome (MECP2 Mutation); Incontinentia Pigmenti; X-Linked Hypophosphatemic Rickets; Aicardi Syndrome; Focal Dermal Hypoplasia (Goltz Syndrome); Fragile X-Associated Tremor/Ataxia Syndrome (FXTAS).

\item[Y-Linked (Holandric) Conditions] Y-Chromosome Infertility (AZF Deletions); SRY Gene Deletions/Mutations; Swyer Syndrome (46,XY Gonadal Dysgenesis); Retinitis Pigmentosa (Y-Linked Form); Hypertrichosis Pinnae Auris.

\item[Mitochondrial (Maternal) Genetic Diseases] Leber Hereditary Optic Neuropathy (LHON); Mitochondrial Encephalomyopathy, Lactic Acidosis, and Stroke-Like Episodes (MELAS); Myoclonic Epilepsy with Ragged Red Fibers (MERRF); Neuropathy, Ataxia, and Retinitis Pigmentosa (NARP); Leigh Syndrome (Subacute Necrotizing Encephalomyelopathy); Kearns–Sayre Syndrome; Pearson Syndrome; Chronic Progressive External Ophthalmoplegia (CPEO).

\item[Chromosomal Aneuploidies (Autosomal)] Trisomy 21 (Down Syndrome); Trisomy 18 (Edwards Syndrome); Trisomy 13 (Patau Syndrome); Trisomy 8 (Warkany Syndrome 2); Trisomy 9 Mosaicism.

\item[Sex Chromosome Aneuploidies] 45,X (Turner Syndrome); 47,XXY (Klinefelter Syndrome); 47,XXX (Triple X Syndrome); 47,XYY (Jacob Syndrome); 48,XXYY Syndrome; 48,XXXY Syndrome; Tetrasomy X.

\item[Microdeletion \& Microduplication Syndromes] 22q11.2 Deletion Syndrome (DiGeorge Syndrome, Velocardiofacial Syndrome); 5p Deletion Syndrome (Cri-du-Chat Syndrome); 4p Deletion Syndrome (Wolf–Hirschhorn Syndrome); 7q11.23 Deletion Syndrome (Williams Syndrome); 17p11.2 Deletion Syndrome (Smith–Magenis Syndrome); 17p11.2 Duplication Syndrome (Potocki–Lupski Syndrome); 1p36 Deletion Syndrome; 15q11–q13 Deletion; Rubinstein–Taybi Syndrome.

\item[Genomic Imprinting Disorders] Prader–Willi Syndrome (Maternal Imprinting / Paternal 15q11–q13 Loss); Angelman Syndrome (Paternal Imprinting / Maternal UBE3A Loss); Beckwith–Wiedemann Syndrome (11p15.5 Alteration); Silver–Russell Syndrome; Pseudohypoparathyroidism (Albright Hereditary Osteodystrophy); Uniparental Disomy Syndromes.

\item[Trinucleotide Repeat Expansion Diseases] Huntington Disease (CAG Repeat); Fragile X Syndrome (CGG Repeat); Myotonic Dystrophy Type 1 (CTG Repeat); Friedreich Ataxia (GAA Repeat); Spinocerebellar Ataxias (SCA Types 1, 2, 3, 6, 7); Dentatorubral-Pallidoluysian Atrophy (DRPLA); Spinal and Bulbar Muscular Atrophy (Kennedy Disease).

\item[Lysosomal Storage Diseases] Gaucher Disease (Glucocerebrosidase Deficiency); Tay–Sachs Disease (Hexosaminidase A Deficiency); Niemann–Pick Disease (Types A, B, C); Fabry Disease (Alpha-Galactosidase A Deficiency); Pompe Disease (Acid Alpha-Glucosidase Deficiency); Krabbe Disease; Metachromatic Leukodystrophy; Mucopolysaccharidoses (Hurler, Scheie, Hunter, Sanfilippo, Morquio, Maroteaux–Lamy, Sly Syndromes); Fucosidosis; Sialidosis.

\item[Connective Tissue Genetic Disorders] Marfan Syndrome (FBN1 Mutation); Ehlers–Danlos Syndromes; Loeys–Dietz Syndrome; Osteogenesis Imperfecta (Brittle Bone Disease); Cutis Laxa; Stickler Syndrome; Pseudoxanthoma Elasticum; Alport Syndrome.

\item[Hereditary Cancer Predisposition Syndromes] Hereditary Breast and Ovarian Cancer Syndrome (BRCA1, BRCA2); Lynch Syndrome (MLH1, MSH2, MSH6, PMS2); Familial Adenomatous Polyposis (APC Mutation); Li–Fraumeni Syndrome (TP53 Mutation); Multiple Endocrine Neoplasia Types 1, 2A, 2B; Von Hippel–Lindau Disease; Cowden Syndrome (PTEN Hamartoma Tumor Syndrome); Peutz–Jeghers Syndrome; Nevoid Basal Cell Carcinoma Syndrome (Gorlin Syndrome); Retinoblastoma (RB1 Mutation); Wilms Tumor Predisposition (WT1).
\end{description}
